\par\noindent\begin{minipage}{\linewidth}
% !TeX root = ../教材题库_全部条目_转写工作稿.tex
% Source: XJ-B1-C03-HW-P085-Y0449
% 第3章 / 3.2 函数的基本性质 / 教材 p.85 / PDF 第92页
% 已对照原页转写题干及公式；待教师复核；未制作答案
% 原题及相邻原解截图（核对用）：../题目原图/XJ-B1-C03-HW-P085-Y0449-01.webp
% 本文件只在转写完成后作为题干使用；OCR 条目尚不是可靠题干。
\begin{bankitem}{习题3.2 第1题}{第3章 / 3.2 函数的基本性质 / 教材 p.85 / PDF 第92页}{XJ-B1-C03-HW-P085-Y0449}{已对照原页转写题干及公式；待教师复核；未制作答案}
如图所示是函数 $f(x)$ 的图象，列出 $f(x)$ 的若干区间，说明它在各区间上的增减性，并指出该函数的最大、最小值点及最值。
\bankfig{C03-HW085-01}{0.68}
\end{bankitem}
\end{minipage}\par\medskip
